\paragraph{window-\(6\) 的 dyadic 方向旗标、仿射截面与奇素分离}
在前述二点纤维方向谱（定理 \ref{thm:terminal-foldbin6-two-point-fiber-direction-spectrum}）、纤维码距谱（定理 \ref{thm:terminal-foldbin6-fiber-hamming-three-valued}）、仿射几何分类（定理 \ref{thm:terminal-foldbin6-fiber-affine-geometry}）与四块 coarse kernel 的 Kirchhoff--Green 算术（推论 \ref{cor:window6-edge-flux-critical-group-cyclic} 与定理 \ref{thm:window6-kirchhoff-green-single-order-padic-rigidity}）之外，window-\(6\) 还呈现一条更细的分层刚性：最小纤维的全部微观几何首先压缩为纯 dyadic 的方向旗标，而真正的奇素算术内容则直到 coarse potential theory 层才首次出现。

\begin{theorem}[boundary 扇区的 dyadic 方向旗标]\label{thm:window6-boundary-dyadic-direction-flag}
\leanverified{paper\_window6\_boundary\_dyadic\_direction\_flag}
记 boundary 三词按
\[
\texttt{100001},\qquad \texttt{100101},\qquad \texttt{101001}
\]
排序。
则其二点纤维方向向量恰为
\[
v(\texttt{101001})=\texttt{100010},\qquad
v(\texttt{100101})=\texttt{100110},\qquad
v(\texttt{100001})=\texttt{111110}.
\]
特别地，上述三向量的支撑形成严格嵌套链
\[
\{1,5\}\subset \{1,4,5\}\subset \{1,2,3,4,5\},
\]
并且它们在 \(\FF_2^6\) 中线性无关。
\end{theorem}

\begin{proof}
由定理 \ref{thm:terminal-foldbin6-two-point-fiber-direction-spectrum} 的精确分类，
\[
v(w)=\texttt{100010}\iff
w\in\{\texttt{000001},\texttt{001001},\texttt{010101},\texttt{101001}\},
\]
\[
v(w)=\texttt{100110}\iff
w\in\{\texttt{000101},\texttt{100101}\},
\]
\[
v(w)=\texttt{111110}\iff
w\in\{\texttt{010001},\texttt{100001}\}.
\]
将其与 boundary 扇区
\(
X_6^{\mathrm{bdry}}=\{\texttt{100001},\texttt{100101},\texttt{101001}\}
\)
相交，即得三条 boundary 纤维的方向向量。
支撑链由三条二进制向量直接读出。
又若
\(
\alpha\,\texttt{100010}+\beta\,\texttt{100110}+\gamma\,\texttt{111110}=0
\)
于 \(\FF_2^6\)，则比较第二坐标得 \(\gamma=0\)，比较第四坐标得 \(\beta=0\)，再比较第五坐标得 \(\alpha=0\)。
故三向量线性无关。证毕。
\end{proof}
\leanverified{paper\_window6\_boundary\_dyadic\_direction\_flag}

\begin{corollary}[boundary 中心的方向实现]\label{cor:window6-boundary-center-direction-realization}
记
\[
B_{\partial}:=\langle \tau_w:\ w\in X_6^{\mathrm{bdry}}\rangle\cong (\ZZ_2)^3
\]
为 boundary sheet-flip 所生成的规范中心子群（推论 \ref{cor:fold-bin-boundary-center-z2}）。
则存在规范单射
\[
\eta:B_{\partial}\hookrightarrow \FF_2^6,
\qquad
\eta(\tau_w):=v(w),
\]
其像恰为
\[
D_{\partial}:=\Span_{\FF_2}\{\texttt{100010},\texttt{100110},\texttt{111110}\}.
\]
因此，window-\(6\) 的 boundary 中心并非仅以抽象 \((\ZZ_2)^3\) 的形式存在，而是在微态超立方体中获得了一个具体的 dyadic 方向基；该基同时携带支撑长度
\[
2<3<5
\]
的嵌套旗标。
\end{corollary}

\begin{proof}
推论 \ref{cor:fold-bin-boundary-center-z2} 已给出
\(
B_{\partial}\cong(\ZZ_2)^3
\)。
由定理 \ref{thm:window6-boundary-dyadic-direction-flag}，三条 boundary 纤维的方向向量线性无关，故上述赋值唯一延拓为一个单射群同态，其像即为所述三维子空间。证毕。
\end{proof}
\leanverified{paper\_window6\_boundary\_center\_direction\_realization}

\begin{theorem}[split-extremal 码距桥]\label{thm:window6-split-extremal-distance-bridge}
\leanverified{paper\_window6\_split\_extremal\_distance\_bridge}
window-\(6\) 的纤维内部最小汉明距离满足三值律
\[
\delta_{\min}(w)\in\{2,3,5\}.
\]
其中极值 \(\delta_{\min}(w)=5\) 恰出现于且仅出现于两条稳定类型
\[
w\in\{\texttt{010001},\texttt{100001}\}.
\]
在 \(B_3\) 字典下，\(\texttt{010001}\) 对应根向量 \(e_2+e_3=(0,1,1)\)，而 \(\texttt{100001}\in X_6^{\mathrm{bdry}}\) 为 boundary 词；并且两条极值纤维共享同一 dyadic 方向
\[
v(\texttt{010001})=v(\texttt{100001})=\texttt{111110}.
\]
因此，window-\(6\) 的最大内部码距并不内生于 cyclic 或 boundary 任一单层，而只在 rigid split 的两侧各出现一次，并由同一方向向量贯穿。
\end{theorem}

\begin{proof}
由定理 \ref{thm:terminal-foldbin6-fiber-hamming-three-valued}，
\(
\delta_{\min}(w)=5
\)
当且仅当
\(
w\in\{\texttt{010001},\texttt{100001}\}
\)。
其中 \(\texttt{100001}\) 是 boundary 词；
而 \(\texttt{010001}\) 在表 \ref{tab:fold6_b3c3_root_dictionary_b3} 中对应根向量 \((0,1,1)=e_2+e_3\)。
再由定理 \ref{thm:terminal-foldbin6-two-point-fiber-direction-spectrum}，
\[
v(w)=\texttt{111110}
\iff
w\in\{\texttt{010001},\texttt{100001}\},
\]
故这两条极值纤维共享同一 dyadic 方向。证毕。
\end{proof}

\begin{corollary}[高重纤维的极值排斥]\label{cor:window6-heavy-fiber-no-max-distance}
记
\[
\mathcal H_{\ge 3}:=\{w\in X_6:\ |F_6^{-1}(w)|\ge 3\}.
\]
则
\[
\#\{w\in \mathcal H_{\ge 3}:\delta_{\min}(w)=2\}=9,\qquad
\#\{w\in \mathcal H_{\ge 3}:\delta_{\min}(w)=3\}=4,\qquad
\#\{w\in \mathcal H_{\ge 3}:\delta_{\min}(w)=5\}=0.
\]
换言之，最大码距 \(5\) 是二点纤维的专属现象；一旦纤维进入三点或四点层，内部几何立即失去这一极端分离。
\end{corollary}

\begin{proof}
由推论 \ref{cor:terminal-foldbin6-fiber-hamming-min-distance-histogram}，
window-\(6\) 的全局码距直方图为
\[
\#\{\delta_{\min}=2\}=13,\qquad
\#\{\delta_{\min}=3\}=6,\qquad
\#\{\delta_{\min}=5\}=2.
\]
另一方面，全部二点纤维共有八条；由定理 \ref{thm:terminal-foldbin6-two-point-fiber-direction-spectrum}，
\(\texttt{100010}\)、\(\texttt{100110}\)、\(\texttt{111110}\) 三类方向分别出现 \(4,2,2\) 次，而它们的汉明权分别为 \(2,3,5\)。
故二点纤维中对应的 \(\delta_{\min}\) 计数为 \((4,2,2)\)。
从全局直方图中扣除，即得 \(\mathcal H_{\ge 3}\) 上的三项计数。证毕。
\end{proof}
\leanverified{paper\_window6\_heavy\_fiber\_no\_max\_distance}

\begin{theorem}[方向谱版线性核障碍]\label{thm:window6-direction-spectrum-linear-kernel-obstruction}
\leanverified{paper\_window6\_direction\_spectrum\_linear\_kernel\_obstruction}
不存在任何线性映射
\[
L:\FF_2^6\to \FF_2^k
\]
使得 \(\mathrm{Fold}^{\mathrm{bin}}_6\) 的纤维划分恰等于 \(L\) 的陪集划分。
更强地，这一不可能性无需诉诸三点或四点纤维；仅由二点纤维方向谱即可强制。
\end{theorem}

\begin{proof}
若存在上述线性映射 \(L\)，则任意大小为 \(2\) 的纤维必为某个仿射陪集
\(
x+\ker L
\)
且 \(\dim\ker L=1\)。
因此所有二点纤维的方向向量都必须等于 \(\ker L\) 中唯一的非零元。
然而定理 \ref{thm:terminal-foldbin6-two-point-fiber-direction-spectrum} 给出三种互异方向
\[
\texttt{100010},\qquad \texttt{100110},\qquad \texttt{111110},
\]
故矛盾。证毕。
\end{proof}

\begin{theorem}[affine \(2\)-flat 的根空间选择]\label{thm:window6-affine-2flat-root-slice-selection}
window-\(6\) 中恰有三条稳定类型的纤维为 affine \(2\)-flat，即
\[
w\in\{\texttt{000000},\texttt{000010},\texttt{101000}\}.
\]
在 \(B_3\) 字典下，它们恰对应于
\[
\texttt{000000}\leftrightarrow (1,1,0),\qquad
\texttt{000010}\leftrightarrow (0,-1,0),\qquad
\texttt{101000}\leftrightarrow (1,-1,0).
\]
因此 affine \(2\)-flat 层并非在根多面体中各向均匀散布，而是精确选出平面 \(z=0\) 上的三点构形
\[
\{e_1+e_2,\ -e_2,\ e_1-e_2\}.
\]
特别地，该层完全位于 cyclic 扇区，并与 boundary 扇区不交。
\leanverified{paper\_window6\_affine\_2flat\_root\_slice\_selection}
\end{theorem}

\begin{proof}
由定理 \ref{thm:terminal-foldbin6-fiber-affine-geometry}，window-\(6\) 中恰有三条 \texttt{affine\_2\_flat} 纤维，对应的稳定类型正是
\[
\texttt{000000},\qquad \texttt{000010},\qquad \texttt{101000}.
\]
再由表 \ref{tab:fold6_b3c3_root_dictionary_b3}，上述三词在 \(B_3\) 规范下分别送到
\((1,1,0)\)、\((0,-1,0)\)、\((1,-1,0)\)。
它们的第三坐标皆为 \(0\)，因而全部落在平面 \(z=0\) 上；同时三词都属于 \(X_6^{\mathrm{cyc}}\)，故与 boundary 扇区无交。证毕。
\end{proof}

\begin{proposition}[\(z=0\) 根截面的手征切割]\label{prop:window6-z0-root-slice-chiral-cut}
\leanverified{paper\_window6\_z0\_root\_slice\_chiral\_cut}
记
\[
\Phi_{z=0}:=\Phi(B_3)\cap\{(x,y,z):z=0\}
=\{(\pm1,\pm1,0),(0,\pm1,0)\}.
\]
则定理 \ref{thm:window6-affine-2flat-root-slice-selection} 选出的 affine \(2\)-flat 根集恰为
\[
\Phi_{z=0}^{\mathrm{aff}}
=\{(1,1,0),(1,-1,0),(0,-1,0)\},
\]
而其补集为
\[
\Phi_{z=0}\setminus \Phi_{z=0}^{\mathrm{aff}}
=\{(-1,1,0),(-1,-1,0),(0,1,0)\}.
\]
换言之，window-\(6\) 在 \(z=0\) 的六根截面上并非做对称一分为二，而是做一条定向切割：保留全部 \(x=1\) 的长根与唯一的负短根 \(-e_2\)，同时排除其镜像三根。
\end{proposition}

\begin{proof}
由定理 \ref{thm:window6-affine-2flat-root-slice-selection}，
\(
\Phi_{z=0}^{\mathrm{aff}}
\)
已确定为前三根。
而 \(B_3\) 根系在平面 \(z=0\) 上的截面恰由
\[
(1,1,0),\ (1,-1,0),\ (-1,1,0),\ (-1,-1,0),\ (0,1,0),\ (0,-1,0)
\]
组成。
取差集即得所述补集。证毕。
\end{proof}

\begin{theorem}[Weyl 对称的图刚性湮灭]\label{thm:window6-weyl-symmetry-graph-annihilation}
\leanverified{paper\_window6\_weyl\_symmetry\_graph\_annihilation}
设 \(\Gamma_6\) 为类型邻接图。
则在由表 \ref{tab:fold6_b3c3_root_dictionary_b3} 固定的 cyclic 根字典下，任何由非平凡 Weyl 元
\(
w\in W(B_3)
\)
在 \(X_6^{\mathrm{cyc}}\) 上诱导的顶点置换，都不能扩张为 \(\Gamma_6\) 的自同构。
等价地，若 \(\sigma\in\Sym(X_6)\) 满足：
\begin{enumerate}
\item \(\sigma|_{X_6^{\mathrm{cyc}}}\) 等于某个 Weyl 元在根字典下诱导的置换；
\item \(\sigma\) 保持 \(\Gamma_6\) 的邻接关系，
\end{enumerate}
则必有 \(\sigma=\mathrm{id}\)。
因此，尽管 cyclic 根层仍携带 \(|W(B_3)|=48\) 的 kinematic 对称，类型邻接层却把它彻底压缩到平凡群。
\end{theorem}

\begin{proof}
由命题 \ref{prop:terminal-gamma6-rigidity}，
\[
\Aut(\Gamma_6)=\{1\}.
\]
因此任何保持 \(\Gamma_6\) 邻接关系的置换都只能是恒等。
若 \(\sigma|_{X_6^{\mathrm{cyc}}}\) 由某个 Weyl 元诱导且 \(\sigma\) 又保持 \(\Gamma_6\) 邻接，则上述刚性强制 \(\sigma=\mathrm{id}\)，从而该 Weyl 作用在类型层也只能平凡。证毕。
\end{proof}

\leanverified{paper\_window6\_local\_determinant\_vs\_coarse\_green\_odd\_primes}
\begin{theorem}[local skeleton 与 coarse Green 的奇素分离]\label{thm:window6-local-determinant-vs-coarse-green-odd-primes}
沿用定理 \ref{thm:terminal-window6-edge-flux-skeleton} 与定理 \ref{thm:window6-kirchhoff-green-single-order-padic-rigidity} 的记号。
则 odd-prime 数据在 window-\(6\) 的 local skeleton 与 coarse potential theory 之间发生严格分层：
\[
\{p\text{ odd prime}:p\mid \det(E)\}=\{3,5,23\},
\]
\[
\{p\text{ odd prime}:p\mid \tau(\mathcal G_E)\}=\{3,571\},
\]
并且对 Green 核
\(
Z=(I-T+\mathbf 1\pi^\top)^{-1}
\)
而言，\(\ZZ[1/6]\)-局部化之后唯一残留的 odd-prime 分母必为 \(571\)，且其 \(571\)-进阶恰为 \(1\)。
特别地，素数 \(23\) 与 \(5\) 在 local skeleton 的整数指纹中可见，却在 coarse Green renormalization 中被精确湮灭；相反，\(571\) 在 \(\det(E)\) 中完全不可见，却在树复杂度与 Green 求逆层首次出现并被不可约保留。
\end{theorem}

\begin{proof}
由 \eqref{eq:window6_edge_flux_skeleton_arithmetic}，
\[
\det(E)=10350=2\cdot 3^2\cdot 5^2\cdot 23,
\]
故 \(\det(E)\) 的奇素支撑为 \(\{3,5,23\}\)。
另一方面，由推论 \ref{cor:window6-edge-flux-critical-group-cyclic}，
\[
\tau(\mathcal G_E)=123336=2^3\cdot 3^3\cdot 571,
\]
故生成树复杂度的奇素支撑为 \(\{3,571\}\)。
最后，由定理 \ref{thm:window6-zeta-cumulant-localization-vs-green-kernel-prime571}，
\(
Z\notin M_4(\ZZ[1/6])
\)；
再由推论 \ref{cor:window6-green-kernel-pstar-valuation-one}，
\(\ZZ[1/6]\)-局部化后的最小公共分母 \(D\) 满足
\[
v_{571}(D)=1.
\]
因此，在 coarse Green 层所有超出 \(\{2,3\}\)-局部化口径的 odd-prime 信息都压缩为唯一的素数 \(571\)。
综上，\(23\) 与 \(5\) 在 coarse Green 层被消去，而 \(571\) 则作为 purely coarse 的新奇素被保留。证毕。
\end{proof}

\begin{conjecture}[dyadic 方向旗标与 odd-prime 分离原理]\label{conj:window6-dyadic-direction-vs-odd-prime-separation}
设偶窗口 \(m\) 同时满足：
\begin{enumerate}
\item \(X_m\) 承认一个 rigid cyclic root dictionary；
\item 最小纤维在 \(\FF_2^m\) 中承认有限且可审计的方向谱；
\item 存在一个有限 coarse kernel，其 Green 矩阵在不可约分母层可被完整审计。
\end{enumerate}
则应有两层互不混同的刚性。
\begin{enumerate}
\item 最小纤维的微观几何全部压缩为一个纯 dyadic 的方向旗标
\[
\mathcal D_m\subset \FF_2^m,
\]
它只记录支撑嵌套、汉明权与 boundary/cyclic 的 split-extremal 现象，而不携带真正的奇素算术信息；
\item 所有真正的奇素算术内容都在 coarse potential theory 层首次出现，并以 Green 核不可约分母的奇素支撑记录下来。
\end{enumerate}
window-\(6\) 正是这一分离原则的首个严格样本：其 dyadic 方向旗标由
\[
\{\texttt{100010},\texttt{100110},\texttt{111110}\}
\]
给出，而 coarse Green 分母的奇素支撑为
\[
\{571\}\quad\text{（相对于 }\ZZ[1/6]\text{-局部化）};
\]
与此同时，local skeleton 的 \(\det(E)\) 只提供 \(\{3,5,23\}\)，其中 \(23\) 与 \(5\) 在 coarse Green renormalization 中被精确消去。
\end{conjecture}
\begin{proposition}[boundary 方向小倍增的最小线性载体]\label{prop:distill-window6-boundary-direction-carrier-extraction}
令
\[
a=\texttt{100010},\qquad b=\texttt{100110},\qquad c=\texttt{111110}\in\FF_2^6
\]
为定理 \ref{thm:window6-boundary-dyadic-direction-flag} 给出的三条 boundary 二点纤维方向，并置
\[
A_{\partial}:=\{0,a,b,c\}\subset\FF_2^6,
\qquad
D_{\partial}:=\Span_{\FF_2}\{a,b,c\}.
\]
则
\[
|A_{\partial}+A_{\partial}|=7<2|A_{\partial}|,
\qquad
A_{\partial}+A_{\partial}+A_{\partial}=D_{\partial}.
\]
并且 \(D_{\partial}\) 是包含 \(A_{\partial}\) 的唯一最小仿射线性载体：若 \(H\leq \FF_2^6\) 为线性子空间且 \(A_{\partial}\subset x+H\)，则 \(D_{\partial}\subset H\)。特别地，任何以这些方向向量实现 boundary sheet-flip 的线性核模型都必须保留三维载体 \(D_{\partial}\)。
\end{proposition}
\begin{proof}
由定理 \ref{thm:window6-boundary-dyadic-direction-flag}，\(a,b,c\) 在 \(\FF_2^6\) 中线性无关。因此
\[
A_{\partial}+A_{\partial}
=\{0,a,b,c,a+b,a+c,b+c\}
\]
且七个元素互异，于是得到第一式。三重和显然包含上述七个元素，并且
\(a+b+c\in A_{\partial}+A_{\partial}+A_{\partial}\)，故三重和包含 \(\Span_{\FF_2}\{a,b,c\}\) 的全部八个元素；反向包含由每一项均在该张成空间中给出。

若 \(A_{\partial}\subset x+H\)，由于 \(0\in A_{\partial}\)，有 \(x\in H\)，从而 \(x+H=H\)。于是 \(a,b,c\in H\)，故 \(D_{\partial}\subset H\)。这说明 boundary 方向片段的二重非增长并不允许压到二维载体；小倍增只能在再乘一次后抽取出精确的三维闭包载体。证毕。
\end{proof}
\begin{theorem}[boundary 中心小倍增闭包的精确承载子抽取]\label{distill:breuillard_green_tao_approximate_groups_and_helfgott_growth-approximate_closure_implies_structured_carrier-boundary-center-sharp-small-doubling-carrier}
沿用推论 \ref{cor:window6-boundary-center-direction-realization} 的记号，令
\[
B_{\partial}=\langle \tau_w:\ w\in X_6^{\mathrm{bdry}}\rangle\cong(\ZZ_2)^3,
\qquad
\eta:B_{\partial}\hookrightarrow \FF_2^6
\]
为 boundary sheet-flip 中心及其 dyadic 方向实现。对非空集合 \(U\subset B_{\partial}\)，记
\[
U^{-1}U:=\{u^{-1}v:\ u,v\in U\}.
\]
若存在 \(u_0\in U\) 使锚定集合
\[
W:=u_0^{-1}U
\]
满足小倍增闭包条件
\[
|U^{-1}U|=|W^2|\le {4\over 3}|U|,
\]
则
\[
H_U:=U^{-1}U=\langle W\rangle
\]
是 \(B_{\partial}\) 的子群，且 \(U\subset u_0H_U\)、
\[
\eta(H_U)=\Span_{\FF_2}\eta(W)\subset D_{\partial}.
\]
更精确地，且仅有下列情形可能发生：
\begin{enumerate}
\item \(W=H_U\)；
\item \(H_U\) 为二维子群，且 \(W=H_U\setminus\{h\}\)，其中 \(0\ne h\in H_U\)；
\item \(H_U=B_{\partial}\)，且 \(6\le |W|\le 7\)。
\end{enumerate}
在所有情形中均有
\[
U^{-1}U=H_U,
\qquad
|H_U|\le {4\over 3}|U|.
\]
因此 boundary 中心上的弱局部闭包可被精确替换为一个 rank 至多 \(3\)、step \(1\) 的 dyadic coset nilprogression 承载子 \(u_0H_U\)。
\end{theorem}

\begin{proof}
经由 \(\eta\) 以及定理 \ref{thm:window6-boundary-dyadic-direction-flag}，可把 \(B_{\partial}\) 视为三维 \(\FF_2\)-向量空间；证明中将群法写作加法，单位元写作 \(0\)。锚定后 \(0\in W\)，且由于 \(B_{\partial}\) 的每个元素均为二阶元，
\[
W+W=U^{-1}U.
\]
令 \(H:=\Span_{\FF_2}(W)\)，并记 \(s:=|W|=|U|\)、\(d:=\dim H\)。显然 \(W+W\subset H\)。下面按 \(d\) 分类。

若 \(d=0\)，则 \(W=\{0\}=H\)，结论成立。若 \(d=1\)，则 \(W\) 含 \(0\) 且张成一条二元直线，故 \(W=H\)，结论亦成立。

若 \(d=2\)，则 \(|H|=4\)。张成二维空间且含 \(0\) 的集合 \(W\) 只能有 \(s=3\) 或 \(s=4\)。当 \(s=4\) 时 \(W=H\)。当 \(s=3\) 时，可写
\[
W=\{0,a,b\},\qquad a,b\ \text{线性无关},
\]
于是
\[
W+W=\{0,a,b,a+b\}=H,
\]
且 \(|W+W|=4=(4/3)|W|\)。这给出第二种情形。

若 \(d=3\)，则 \(H=B_{\partial}\)，且 \(s\ge 4\)。当 \(s=4\) 时，若
\[
W=\{0,a,b,c\}
\]
并且 \(a,b,c\) 线性无关，则
\[
0,a,b,c,a+b,a+c,b+c
\]
为 \(W+W\) 中七个互异元素，故
\[
|W+W|\ge 7>{4\over 3}\cdot 4,
\]
与小倍增闭包条件矛盾。当 \(s=5\) 时，由 \(|W|> |B_{\partial}|/2=4\)，对任意 \(x\in B_{\partial}\)，集合 \(W\) 与 \(x+W\) 必相交，故 \(x\in W+W\)。于是 \(W+W=B_{\partial}\)，从而
\[
|W+W|=8>{4\over 3}\cdot 5,
\]
仍矛盾。故三维张成情形下只能有 \(s=6,7,8\)。其中 \(s=8\) 已包含在 \(W=H\) 的闭合情形内；\(s=6,7\) 给出第三种情形。并且同样由 \(|W|>4\) 可知 \(W+W=B_{\partial}=H\)。

以上分类同时证明了 \(U^{-1}U=W+W=H=\langle W\rangle\)，以及
\[
|H|=|U^{-1}U|\le {4\over 3}|U|.
\]
若改取另一锚点 \(u_1\in U\)，则 \(u_1^{-1}u_0\in U^{-1}U=H\)，因此 \(u_1^{-1}U\) 仍张成同一子群 \(H\)，故 \(H_U\) 与锚点无关。

最后，取 \(H_U\) 的任意 \(\FF_2\)-基 \(h_1,\dots,h_r\)，其中 \(r\le 3\)。则
\[
H_U=\{h_1^{\epsilon_1}\cdots h_r^{\epsilon_r}:\epsilon_i\in\{0,1\}\}
\]
是阿贝尔 step \(1\) nilprogression，而 \(u_0H_U\) 是相应 coset nilprogression。其方向像正是
\(\eta(H_U)=\Span_{\FF_2}\eta(W)\subset D_{\partial}\)。这正是在 boundary dyadic 中心内的小倍增闭包抽取：弱闭包集合 \(U^{-1}U\) 被强制等同于一个显式结构承载子。证毕。
\end{proof}

\endinput
